\hypertarget{pell-harmonic-a-harmonic-scale-transition-in-the-pell-cat-map}{%
\section{A harmonic-scale transition in the Pell cat map}\label{pell-harmonic-a-harmonic-scale-transition-in-the-pell-cat-map}}

\textbf{Status: project proof, with explicit scope.} This note derives a family of
global observable laws in one fixed arithmetic quantum map. It does not prove
microscopic BGS statistics. The limiting graphon moment construction is prior
work; the claim here is its realization by growing harmonics of the Pell cat
map. Priority for that realization has not been established.

\hypertarget{pell-harmonic-statement-and-normalization}{%
\subsection{Statement and normalization}\label{pell-harmonic-statement-and-normalization}}

Let \(K\) tend to infinity through odd primes. Write
\[
 \alpha=1+\sqrt2,\quad \lambda=\alpha^2,\qquad
 N+F\sqrt2=\alpha^K,\qquad T=\begin{pmatrix}3&2\\4&3\end{pmatrix}.
\]
Use the parity-preserving quantization \(U_0\) of this fixed matrix, normalized
by \(U_0^K=I\), as in the preceding Pell block proof.
The odd parity space has dimension \(D=(N-1)/2\). Its \(K\) eigenspaces
\(E_\eta\) have equal dimension \(n=(N-1)/(2K)\); let \(P_\eta\) be their projections.

For a positive integer \(a\) coprime to \(N\), take the actual physical harmonic
\[
 B_a=W(a,0)+W(-a,0),\qquad
 (B_a u)(x)=2\cos(2\pi ax/N)u(x).
\]
Here \(W\) has the modular-centered convention of the Pell block proof.
Define
\[
 X_{K,a,\eta}=\sqrt{K/2}\,(P_\eta B_aP_\eta)|_{E_\eta},
 \qquad \mu_{m,K,\eta}(a)=n^{-1}\operatorname{Tr}X_{K,a,\eta}^m.       \tag{H1}
\]
Let \(1\le L_K\le\lfloor\sqrt N\rfloor\) be integers with
\[
 \frac{\log L_K}{\log N}\longrightarrow\beta\in[0,1/2],
 \qquad \mathcal A_K=\{1\le a\le L_K:\gcd(a,N)=1\}.                 \tag{H2}
\]
Probability and \(L^2\) below use the uniform distribution on \(\mathcal A_K\).
No randomness is added to the map or to its quantization.

Put \(\rho=1-2\beta\) and define on the unit circle
\[
 w_\beta(x,y)=\mathbf1\{\operatorname{dist}_{\mathbb R/\mathbb Z}(x,y)
                                  <\rho/2\}.                    \tag{H3}
\]
Endpoint values on the distance boundary do not affect integrals. For a pair
partition \(\pi\) of \(2r\) positions, its crossing graph \(G_\pi\) has one vertex
per pair and an edge precisely when the two pairs interlace. Define
\[
 t(G,w)=\int_{[0,1]^{V(G)}}\prod_{\{i,j\}\in E(G)}w(x_i,x_j)\,dx .
\]
\[
 M_{2r}(\beta)=\sum_{\pi\in\mathcal P_2(2r)}t(G_\pi,w_\beta),
 \quad M_{2r+1}(\beta)=0,\quad M_0=1.                            \tag{H4}
\]

\textbf{Theorem 1.} For every fixed nonnegative integer \(m\),
\[
 \left\|\sup_\eta|\mu_{m,K,\eta}-M_m(\beta)|\right\|_{L^2(\mathcal A_K)}
       \longrightarrow0.                                     \tag{H5}
\]
There is a unique probability measure \(\nu_\beta\) with moments (H4).
The empirical spectral measures of (H1) converge weakly to \(\nu_\beta\) in
probability, uniformly over all \(K\) sectors. The averaged measures converge
weakly too. At \(\beta=0\) the limit is the standard Gaussian; at \(\beta=1/2\)
it is the variance-one semicircle.

For every \(0\le\beta<1/2\), \(\nu_\beta\) has unbounded support. Thus the
compactly supported endpoint is a genuine change of support, despite weak
continuity of the family.

The intermediate laws are not the ordinary \(q\)-Gaussian laws with
\(q=\rho\). In particular,
\[
 M_2=1,\qquad M_4=2+\rho=3-2\beta,
\]
\[
 M_6=5+6\rho+\frac{15}{4}\rho^2+
                         \left(\frac32\rho-1\right)_+^2.        \tag{H6}
\]
Thus (H5) is an all-fixed-order statement, not an interpolation inferred
from the fourth moment.

\hypertarget{pell-harmonic-why-the-harmonic-enters-quadratically}{%
\subsection{Why the harmonic enters quadratically}\label{pell-harmonic-why-the-harmonic-enters-quadratically}}

Set \(s=F^{-1}\pmod N\), \(h=3+2s\), and \(q_0=-F\). In the diagonal coordinates
of the Pell block proof, the fixed conjugation taking \(B_1\) to
\(W(q_0,1)+W(q_0,1)^*\) takes \(B_a\) to
\[
 W(aq_0,a)+W(aq_0,a)^*.
\]
The dilation \(S_u f(x)=f(u^{-1}x)\) satisfies
\[
 S_uW(b,c)S_u^{-1}=W(b/u,uc).
\]
Taking \(u=a^{-1}\) converts this to \(W(a^2q_0,1)+W(a^2q_0,1)^*\).
This dilation commutes with \(S_h\) and parity, so it preserves every sector.
Consequently (H1) is unitarily equivalent to the previously defined block
\(X_{q,\eta}\) with \(q=a^2q_0\), up to the fixed initial permutation of sector
labels. It is not the block with \(q=aq_0\).

For a residue \(d\) modulo \(K\), let \(\delta(d)\) be its signed representative
in \([-(K-1)/2,(K-1)/2]\). Extend the even Pell sequence by
\(P_{-2d}=-P_{2d}\). The exact identity is
\[
 q_0(h^d-h^{-d})\equiv-2P_{2\delta(d)}\pmod N,\qquad
 P_{2d}=\frac{\lambda^d-\lambda^{-d}}{2\sqrt2}.                  \tag{H7}
\]
Thus a crossing contributes a phase quadratic in \(a\).

\hypertarget{pell-harmonic-the-unit-restriction-has-vanishing-cost}{%
\subsection{The unit restriction has vanishing cost}\label{pell-harmonic-the-unit-restriction-has-vanishing-cost}}

Every prime \(\ell\mid N\) admits an element \(h\) of exact order \(K\):
\(h^K=1\) and \(h-1\) is a unit, so reduction modulo \(\ell\) cannot have order
one. Therefore \(\ell\equiv1\pmod{2K}\) and \(\ell\ge2K+1\). Counting distinct
prime divisors gives
\[
 \frac{\#([1,L_K]\setminus\mathcal A_K)}{L_K}
 \le\sum_{\ell\mid N}\frac1\ell
 \le\frac{\log N}{(2K+1)\log(2K+1)}
 =O(1/\log K).                                               \tag{H8}
\]
The first inequality uses \(\lfloor L_K/\ell\rfloor\le L_K/\ell\) and
is valid even for bounded \(L_K\). We will perform phase averaging over the
whole interval and condition afterwards. All averaged phase expressions
are bounded in modulus; no unbounded-moment conditioning is implicit.

\hypertarget{pell-harmonic-reduction-to-pairings-uniformly-in-the-harmonic}{%
\subsection{Reduction to pairings, uniformly in the harmonic}\label{pell-harmonic-reduction-to-pairings-uniformly-in-the-harmonic}}

The all-order trace identity and largest-gap lifting argument of the Pell
block proof apply uniformly to every unit \(q\), hence to \(q=a^2q_0\).
For completeness, a central signed \(m\)-word gives both reciprocal
congruences. After cutting at its largest cyclic gap, its signed polynomial
\(P\) has degree at most \(K-\lceil K/m\rceil\) and coefficient norm at most \(m\).
Its remainder modulo \(Z^2-6Z+1\) is divisible coefficientwise by \(N\).
A nonzero quadratic norm would therefore satisfy
\[
 N^2\le|\operatorname{Norm}P(\lambda)|
       \le m^2\lambda^{K-\lceil K/m\rceil},
\]
which is impossible for sufficiently large \(K\) at fixed \(m\).
Thus the relation lifts to \(P(\lambda)=0\).

A gap greater than \(2\log_\lambda m\) splits this relation into separately
vanishing pieces: if the high piece were nonzero, its value at \(\lambda\)
would have modulus at least \(1/m\) by its integral quadratic norm, while
the lower part cannot cancel it across such a gap. Each nonzero piece
has bounded span and at least three signed occurrences. Reinsert cancelled
opposite pairs as two-occurrence pieces. A nonbalanced \(2r\)-word then has
at most \(r-1\) translation parameters; after pinning one label there are
\(O_r(K^{r-2})\) such words. The normalized contribution is \(O_r(K^{-1})\),
independently of their phases. For odd \(m\), the equality \(P(\lambda)=0\)
implies \(4\mid P(1)\), impossible for an odd number of signs.

Balanced words are counted by pair partitions, opposite orientations
\(\sigma_i\in\{-1,1\}\), and labels \(e_i\in\mathbb Z/K\mathbb Z\).
Duplicated labels contribute \(O_r(K^{-1})\) after normalization.
Orient pairs by their first endpoints. For a crossing graph with \(r\)
vertices, set
\[
 Z_{\pi,\sigma,K}(a)=K^{-(r-1)}
 \sum_{\substack{e_1=0\\e_2,\ldots,e_r\bmod K}}
       e\!\left(a^2 u_{\pi,\sigma,K}(e)\right),\qquad e(t)=e^{2\pi it},
\]
\[
 u_{\pi,\sigma,K}(e)=
 -\frac2N\sum_{\substack{i<j\\\{i,j\}\in E(G_\pi)}}
       \sigma_i\sigma_j P_{2\delta(e_j-e_i)}.                    \tag{H9}
\]
The exact Weyl multiplication formula gives
\[
 \mu_{2r,K,\eta}(a)
 =\frac{N}{N-1}\sum_{\pi}2^{-r}\sum_\sigma Z_{\pi,\sigma,K}(a)
       +O_r(K^{-1})+O_r(K^{r+1}/N).                            \tag{H10}
\]
The errors are uniform over all unit \(a\) and all sectors. Odd moments
are uniformly \(O_m(K^{m/2+1}/N)\). The common leading expression in
(H10), rather than a union bound over \(K\) sectors, will give (H5).

\hypertarget{pell-harmonic-an-elementary-quadratic-sum-estimate}{%
\subsection{An elementary quadratic-sum estimate}\label{pell-harmonic-an-elementary-quadratic-sum-estimate}}

\textbf{Lemma 2.} Suppose \(u=u_K\) is real and
\[
 0<|u|\le N^{-c},\qquad L_K\sqrt{|u|}\ge N^c                 \tag{H11}
\]
for some fixed \(c>0\). Then
\[
 L_K^{-1}\sum_{a=1}^{L_K}e(ua^2)\longrightarrow0
\]
uniformly under (H11).

\textbf{Proof.} For \(z_a=e(ua^2)\) on \(1\le a\le L\), zero elsewhere, Cauchy--Schwarz
applied to \(H\sum z_a=\sum_b\sum_{j=0}^{H-1}z_{b-j}\) gives, for \(H\le L\),
\[
 \left|L^{-1}\sum z_a\right|^2
 \le\frac2H+\frac4{HL}\sum_{j=1}^{H-1}
                   \min\!\left(L,\frac1{2\|2uj\|_{\mathbb R/\mathbb Z}}\right).
                                                               \tag{H12}
\]
Take \(H=\lfloor(4\sqrt{|u|})^{-1}\rfloor\). Eventually \(2\le H\le L\) and
\(\|2uj\|=2|u|j\) for \(1\le j<H\). The geometric-sum bound and the harmonic
sum therefore imply
\[
 \left|L^{-1}\sum z_a\right|^2
 \ll \sqrt{|u|}+
      \frac{\log(2/|u|)}{L\sqrt{|u|}}.                         \tag{H13}
\]
In our application \(|u|\ge N^{-C}\) for fixed \(C\), so (H11)--(H13) tend
to zero. In fact (H11) and \(L\le\sqrt N\) already imply such a polynomial
lower bound. This proves the lemma without assuming independence or
irrationality of \(u\). Rational phases are included.

\hypertarget{pell-harmonic-the-largest-crossing-distance-controls-the-phase}{%
\subsection{The largest crossing distance controls the phase}\label{pell-harmonic-the-largest-crossing-distance-controls-the-phase}}

\textbf{Lemma 3.} For every fixed pair partition and orientation,
\[
 \|Z_{\pi,\sigma,K}-t(G_\pi,w_\beta)\|_{L^2([1,L_K])}
       \longrightarrow0.                                    \tag{H14}
\]

\textbf{Proof.} If there are no edges, both expressions are one.
Otherwise regard \(x_i=e_i/K\) as points on the circle, with \(x_1=0\),
and let \(d_{ij}=\operatorname{dist}(x_i,x_j)\) for graph edges.
Fix a small \(\varepsilon>0\). Discard configurations for which an edge
distance is within \(\varepsilon\) of \(0\), \(1/2\), or \(\rho/2\), or two
distinct edge distances differ by at most \(\varepsilon\).
Their normalized number is \(O_r(\varepsilon+K^{-1})\).

Here is the geometric justification, including dependent edges. Each circle
distance is piecewise affine in the labels. On each of finitely many cells,
distance equality for two distinct edges is a nontrivial affine equation:
their endpoint coefficient vectors differ, even when the edges share a
vertex or one vertex is fixed at zero. Its \(\varepsilon\)-neighborhood has
volume \(O_r(\varepsilon)\) and grid count \(O_r(\varepsilon K^{r-1}+K^{r-2})\).
The same argument handles the three single-distance boundaries. No assertion
of independent edge distances is made.

On the remaining configurations, let \(d_*=\max d_{ij}\). Its edge is unique
and is separated from every other edge by at least \(\varepsilon\).
Formula (H7) implies that the associated Pell term dominates the sum
in (H9), regardless of orientations. More precisely,
\[
 |u_{\pi,\sigma,K}(e)|\asymp_r
       \lambda^{Kd_*-K/2}\asymp_r N^{2d_*-1},                 \tag{H15}
\]
uniformly on these configurations for sufficiently large \(K\).
The ratio of all other terms to the largest is
\(O_r(\lambda^{-\varepsilon K})\). The exclusion near \(1/2\) ensures
\(|u|\ll_r N^{-2\varepsilon}\): there is no hidden reduction near a nonzero
integer in applying Lemma 2.

If \(d_*<\rho/2-\varepsilon\), then
\(L_K^2|u|\to0\) uniformly, so \(e(a^2u)\to1\) uniformly for \(a\le L_K\).
If \(d_*>\rho/2+\varepsilon\), then
\(L_K\sqrt{|u|}\ge N^{\varepsilon/2}\) eventually; (H13) makes its
average over \(a\) tend to zero. Taking Riemann sums, then letting
\(\varepsilon\downarrow0\), proves
\[
 \mathbb E_a Z_{\pi,\sigma,K}\longrightarrow t(G_\pi,w_\beta).
                                                               \tag{H16}
\]
Pinning \(x_1=0\) does not change this integral, by translation invariance.

For the second moment use two independent label configurations \(e,f\)
and expand
\[
 \mathbb E_a|Z_{\pi,\sigma,K}(a)|^2
 =\mathbb E_{e,f}\mathbb E_a
           e\!\left(a^2[u_{\pi,\sigma,K}(e)-u_{\pi,\sigma,K}(f)]\right).
                                                               \tag{H17}
\]
In addition discard ties between an edge distance in one copy and an edge
distance in the other. They again have measure and grid proportion
\(O_r(\varepsilon+K^{-1})\), since the two copies have independent free
coordinates. Outside this set there is a unique largest edge among both
copies. Its coefficient is \(1\) or \(-1\) and cannot cancel.
The same low/high dichotomy shows that only configurations where \textbf{all}
edges of both copies have distance below \(\rho/2\) survive. Their limiting
volume is \(t(G_\pi,w_\beta)^2\). Together with (H16), this proves (H14).
At \(\beta=0\) only the low region survives; at \(\beta=1/2\) only the high
region survives for a graph with an edge. Thus both endpoints are included.

\hypertarget{pell-harmonic-completion-of-the-spectral-limit}{%
\subsection{Completion of the spectral limit}\label{pell-harmonic-completion-of-the-spectral-limit}}

Since \(|Z_{\pi,\sigma,K}|\le1\), (H8) transfers (H14) to \(\mathcal A_K\).
Sum over the finitely many pairings and orientations in (H10). Its uniform
remainder proves (H5), including the supremum over sectors.

To justify the probability measure and weak convergence, note
\[
 0\le M_{2r}(\beta)\le(2r-1)!!.
\]
The averaged empirical measures have bounded second moments and hence
are tight. Their uniformly bounded higher even moments, supplied by
(H10), give uniform integrability for each lower moment. Any weak
subsequential limit therefore has moments (H4), proving existence.
The displayed Gaussian moment bound gives
\(\int e^{t x^2}\,d\nu_\beta(x)<\infty\) for \(0<t<1/2\) by monotone expansion.
It implies moment determinacy (expand the characteristic function under
this domination). Thus the measure is unique.

For the random empirical laws, from any subsequence choose a further one
on which the suprema in (H5) tend to zero almost surely for every integer
\(m\), by diagonal extraction. Tightness and higher-moment uniform
integrability then give weak convergence for every choice of sector,
including choices depending on \(a\). A contradiction using a maximizing
sector proves uniform weak convergence in probability, equivalently
convergence of the supremum of bounded-Lipschitz distances. Boundedness of
that distance also gives the averaged statement. At the endpoints (H4)
counts respectively all pairings and only noncrossing pairings.

\textbf{Corollary 4 (support transition and continuity).} For \(\rho>0\), choose
one circle interval \(I\) of length \(b=\rho/2\). On \(I\times I\) the kernel is
one almost everywhere. For every crossing graph on \(r\) vertices,
\(t(G,w_\beta)\ge b^r\) by restricting all labels to \(I\). Consequently
\[
 b^r(2r-1)!!\le M_{2r}(\beta)\le(2r-1)!!.                    \tag{H22}
\]
The \(2r\)-th root of the lower bound tends to infinity: the last
\(\lfloor r/2\rfloor\) factors of \((2r-1)!!\) are at least \(r\).
A probability measure supported on \([-R,R]\) would instead have all these
roots bounded by \(R\). Hence the support is unbounded for \(\beta<1/2\).
For \(\beta=1/2\) only noncrossing pairings survive, giving support \([-2,2]\).

For each fixed graph, its density is continuous in \(\beta\) by dominated
convergence, since each distance boundary has measure zero. Thus every
fixed moment is continuous, uniformly bounded by the Gaussian moments.
Tightness, uniform integrability, and moment determinacy prove weak
continuity of \(\beta\mapsto\nu_\beta\), including at the compact endpoint.
This is a statement about the limiting measures, not finite-\(N\) extreme
eigenvalue rigidity.

\hypertarget{pell-harmonic-sixth-moment-and-the-failure-of-an-independent-crossing-model}{%
\subsection{Sixth moment and the failure of an independent-crossing model}\label{pell-harmonic-sixth-moment-and-the-failure-of-an-independent-crossing-model}}

The fifteen pairings of six positions have respectively \(0,1,2,3\) crossings
with multiplicities \(5,6,3,1\). The two-edge graphs are trees, whose
densities are \(\rho^2\) because \(w_\beta\) has constant degree \(\rho\).
The last graph is a triangle. Put \(b=\rho/2\). Fix its first vertex at zero:
\[
 t(K_3,w_\beta)
 =\int_{-b}^b\int_{-b}^b
        \mathbf1\{\operatorname{dist}(x,y)<b\}\,dx\,dy
 =3b^2+(3b-1)_+^2.                                        \tag{H18}
\]
Indeed the central region \(|x-y|<b\) has area \(3b^2\); when \(3b>1\),
the two wrapped corner triangles add \((3b-1)^2\). This proves (H6).
An independent-crossing model would instead assign \(\rho^3\) to the
triangle. The difference is
\[
 t(K_3,w_\beta)-\rho^3=
 \begin{cases}
  \rho^2(3/4-\rho),&0\le\rho\le2/3,\\
  (1-\rho)^3,&2/3\le\rho\le1.
 \end{cases}                                               \tag{H19}
\]
It is strictly positive for \(0<\rho<1\). For example, at \(\beta=1/4\),
\(M_4=5/2\), \(M_6=143/16\), whereas that \(q\)-Gaussian model gives \(71/8\).

\hypertarget{pell-harmonic-weighted-quantum-fields-and-a-noncommutativity-diagnostic}{%
\subsection{Weighted quantum fields and a noncommutativity diagnostic}\label{pell-harmonic-weighted-quantum-fields-and-a-noncommutativity-diagnostic}}

There is also an ordered mixed-moment version on the full odd space.
For a bounded real Riemann-integrable profile \(f\) on \([0,1)\), set
\[
 Y_{K,a}(f)=\frac1{\sqrt{2K}}\sum_{j=0}^{K-1}
                 f(j/K)U_0^jB_aU_0^{-j},\qquad
 \tau_-(A)=D^{-1}\operatorname{Tr}_-A .
\]
For fixed profiles \(f_1,\ldots,f_{2r}\), their ordered mixed trace converges
in \(L^2(\mathcal A_K)\) to
\[
 \sum_{\pi\in\mathcal P_2(2r)}
 \int_{[0,1]^r}
   \prod_{\{u,v\}\in\pi} f_u(x_{\{u,v\}})f_v(x_{\{u,v\}})
   \prod_{\{i,j\}\in E(G_\pi)}w_\beta(x_i,x_j)\,dx.             \tag{H20}
\]
Odd mixed traces tend to zero. In proving this version, use all \(r\) labels
without pinning one; keep the profiles inside the sums. The same discarded
sets and quadratic-phase bounds apply with a fixed supremum-norm factor.
The product profiles and indicator kernels are Riemann integrable, so the
remaining grid sums converge. The noncentral full-space trace error is
\(O_{r,f}(K^r/N)\) and the unbalanced/repeated-label error is \(O_{r,f}(K^{-1})\).
This proves (H20); it is not an assumption of finite quantum independence.

At \(\beta=0\), (H20) is the commuting Gaussian field with covariance
\(\int fg\); at \(\beta=1/2\) it is the free semicircular field with that
covariance. These descriptions concern tracial moments.
For \(f=\sqrt2\,\mathbf1_{[0,1/2)}\) and
\(g=\sqrt2\,\mathbf1_{[1/2,1)}\), (H20) gives
\[
 \tau_-\!\left([Y_{K,a}(f),Y_{K,a}(g)]^*
                    [Y_{K,a}(f),Y_{K,a}(g)]\right)
 \longrightarrow 2(1-\rho^2)=8\beta(1-\beta)                 \tag{H21}
\]
in \(L^2(\mathcal A_K)\). To check the coefficient, for any disjointly supported
real profiles the limit is
\(2\int f(x)^2g(y)^2[1-w_\beta(x,y)]\,dx\,dy\).
The area of allowed pairs across the two half-circles is \(b^2\).
This gives (H21), and shows why the limiting intermediate fields cannot
be treated as commuting random variables.

\hypertarget{pell-harmonic-prior-work-verification-and-the-remaining-bgs-gap}{%
\subsection{Prior work, verification, and the remaining BGS gap}\label{pell-harmonic-prior-work-verification-and-the-remaining-bgs-gap}}

Cébron, Oliveira Santos, and Youssef,
\emph{Graphon-Theoretic Approach to Central Limit Theorems for
\(\epsilon\)-Independence}, \href{https://arxiv.org/html/2411.13062v2}{arXiv:2411.13062v2},
Theorem 1.1 and Section 2,
already give the graphon moment family (H4), including Gaussian and
semicircular endpoints. Their weighted and multivariate extensions also
precede this note. The present Weyl orbit is not assumed to satisfy their
finite independence hypotheses. The additional argument here is the
arithmetic harmonic realization (H7)--(H17). Its priority is unverified.

That paper's Proposition 2.6 assumes \(\|w\|_\infty<1\) to obtain a density.
For \(0\le\beta<1/2\), our kernel has essential supremum one, so that
proposition does not automatically prove absolute continuity here.
No density or finite-\(N\) spectral-edge assertion is being imported; the
limiting support statement follows instead from (H22).

The dedicated verifier checks the physical harmonic/squared-phase
identification against finite matrices, exact Pell and oriented-word
identities, and the low-order crossing geometry. The companion numerical
experiment samples harmonics with integer modular phases; sampling and
floating-point output are not an asymptotic proof. Source hashes and actual
run ranges belong to the dated review record.

The classical map is fixed, but the observable frequency and its distribution
vary with \(N\). These are global spectra of compressed observables and ordered
moments of orbit sums. No local spacing law for \(U_0\), perturbed-propagator
universality, Heisenberg-time trace comparison, or absence of arithmetic
symmetries follows. BGS remains open.
